\subsection{道路}

定义 1. --- 设 X 为拓扑空间。X 中的道路是指从 I 到 X 中的任一连续映射 c。点 \(c(0)\) 称为道路 c 的起点，点 \(c(1)\) 称为道路 c 的终点。X 中的圈是指起点与终点相等的 X 中的道路。

设 \(x\) 与 \(y\) 为 X 的点。称 X 中的道路 \(c\) 连接 \(x\) 与 \(y\)，如果 \(x\) 是它的起点且 \(y\) 是它的终点。

定义 2. --- 设 X 为拓扑空间。若 \(c(1) = d(0)\)，则称 X 中的道路 c 与 d 是可拼接的。这时把由下列公式定义的道路 \(c*d\) 称为 c 与 d 的拼接道路：

\[
(c * d) (t) = \left\{ \begin{array}{l l} c (2 t) & \text { for } 0 \leqslant t \leqslant 1 / 2, \\ d (2 t - 1) & \text { for } 1 / 2 \leqslant t \leqslant 1. \end{array} \right.\tag{1}
\]

它的起点是 c 的起点，它的终点是 d 的终点。

设 c 为 X 中的道路；把 c 的逆道路记作 \(\overline{c}\)，它是由 \(\overline{c}(t) = c(1 - t)\)（对 \(t \in I\)）定义的道路。

若 c 与 d 是 X 中两条可拼接的道路，则 \(\overline{d}\) 与 \(\overline{c}\) 也可拼接，且有 \(\overline{c*d} = \overline{d} * \overline{c}\)。对 X 中每条道路 c，有 \(\overline{\overline{c}} = c\)。

注记。--- 1) 设 X 为拓扑空间，P 为缩化为一点的拓扑空间。借助投影 pr₂ 把 P × I 与 I 等同。从 P 到 X 的映射（必然连续）之间的同伦组成的集合 C(P × I; X) 这时等同于 X 中的道路组成的集合 C(I; X)。连接以 x 为像的常值映射与以 y 为像的常值映射的同伦，对应于以 x 为起点、y 为终点的道路。上述等同与拼接及取逆道路的概念相容（参见 III, p. 230）。

\begin{enumerate}
\def\labelenumi{\arabic{enumi})}
\setcounter{enumi}{1}
\tightlist
\item
  设 X 与 Y 为拓扑空间。典范映射
\end{enumerate}

\[
\mathcal {C} (\mathrm{X} \times \mathbf {I}; \mathrm{Y}) \rightarrow \mathcal {C} (\mathbf {I}; \mathcal {C} _ {c} (\mathrm{X}; \mathrm{Y}))
\]

把每个同伦 \(\sigma\colon X \times I \to Y\)（它连接从 X 到 Y 的两个映射 \(f\) 与 \(g\)）对应到一条以 \(f\) 为起点、\(g\) 为终点的、空间 \(\mathcal{C}_c(X;Y)\) 中的道路。若 \(X\) 是局部紧空间，这个典范映射是双射（TG, X, p. 28, 定理 3）。

设 X 为拓扑空间。把 X 的道路空间记作 \(\Lambda(\mathrm{X})\)，它就是拓扑空间 \(\mathcal{C}_{c}(\mathbf{I};\mathrm{X})\)，其元素是 X 中的道路，其拓扑是紧收敛拓扑（参见 TG, X, p. 27）。若 x 是 X 的一点，用 \(\Lambda_{x}(\mathrm{X})\) 表示 \(\Lambda(\mathrm{X})\) 中由以 x 为起点的道路组成的子空间。若 y 是 X 的另一点，用 \(\Lambda_{x,y}(\mathrm{X})\) 表示 \(\Lambda(\mathrm{X})\) 中由以 x 为起点、y 为终点的道路组成的子空间。

命题 1. --- 设 X 与 Z 为拓扑空间。为使映射 \(\varphi\)（从 Z 到道路空间 \(\mathcal{C}_c(\mathbf{I};\mathrm{X})\) 中）连续，必须而且只需由 \((z,t)\mapsto \varphi (z)(t)\) 给出的映射是从 \(\mathbf{Z}\times \mathbf{I}\) 到 X 中的连续映射。特别地，映射 \((c,t)\mapsto c(t)\)（从 \(\mathcal{C}_c(\mathbf{I};\mathrm{X})\times \mathbf{I}\) 到 X 中）是连续的。

由于拓扑空间 I 是局部紧的，本命题由 TG, X, p. 28, 定理 3 得出。

把映射 \(c \mapsto c(0)\)（从 \(\Lambda(\mathrm{X})\) 到 X 中）记作 o，称为起点映射；把映射 \(c \mapsto c(1)\)（从 \(\Lambda(\mathrm{X})\) 到 X 中）记作 e，称为终点映射。映射 o 与 e 是连续的（TG, X, p. 27, 注 1）。X 中可拼接道路的对就是 X-空间 \((\Lambda(\mathrm{X}), e)\) 与 \((\Lambda(\mathrm{X}), o)\) 的纤维积的元素。

命题 2. --- 设 X 为拓扑空间。把道路 c 对应到逆道路 \(\overline{c}\) 的映射是 \(\Lambda(X)\) 到其自身上的同胚。道路的拼接 \((c,d)\mapsto c*d\) 是空间 \((\Lambda(X),e)\times_{X}(\Lambda(X),o)\) 到 \(\Lambda(X)\) 上的同胚。

映射 \(t \mapsto 1 - t\) 是空间 I 到其自身上的同胚；第一个断言由此得出。

用 C 表示纤维积 \((\Lambda(\mathrm{X}),e)\times_{\mathrm{X}}(\Lambda(\mathrm{X}),o)\)。用 \(\gamma:C\to\Lambda(X)\) 表示映射 \((c,d)\mapsto c*d\)，用 \(\delta:C\times I\to X\) 表示映射 \(((c,d),t)\mapsto(c*d)(t)\)。由公式 (1) 及命题 1，映射 \(\delta\) 在 \(C\times[0,\frac{1}{2}]\) 与 \(C\times[\frac{1}{2},1]\) 上的限制是连续的。因此映射 \(\delta\) 是连续的（TG, I, p. 19, 命题 4），映射 \(\gamma\) 亦然（命题 1）。我们有 \(c(t)=(c*d)(\frac{t}{2})\) 与 \(d(t)=(c*d)(\frac{1+t}{2})\)，因此 \(\gamma\) 是单射。

设 \(g\) 为 \(\mathbf{X}\) 中的道路；对 \(t \in \mathbf{I}\)，令

\[
c _ {g} (t) = g \bigl (\frac {t}{2} \bigr) \quad \text{且} \quad d _ {g} (t) = g \bigl (\frac {1 + t}{2} \bigr).
\]

这样定义的映射 \(c_{g}\) 与 \(d_{g}\) 是 X 中可拼接的道路，且有 \(c_{g} * d_{g} = g\)。此外，映射 \(g \mapsto c_{g}\) 与 \(g \mapsto d_{g}\) 是空间 \(\Lambda(\mathrm{X})\) 到其自身的连续映射（命题 1）。由此推出映射 \(\gamma\) 是同胚。

推论. --- 设 X 为拓扑空间，x、y、z 为 X 的点。映射 \(c \mapsto \overline{c}\)（从 \(\Lambda_{x,y}(X)\) 到 \(\Lambda_{y,x}(X)\)）与映射 \((c,d) \mapsto c * d\)（从 \(\Lambda_{x,y}(X) \times \Lambda_{y,z}(X)\) 到 \(\Lambda_{x,z}(X)\)）都是连续的。

